
\documentclass[11pt,a4paper]{article}

% Page layout
\usepackage[margin=1in]{geometry}

% Mathematics
\usepackage{amsmath,amssymb,amsthm,mathtools}
\usepackage{fontspec}
\usepackage{hyperref}

% Bibliography
\usepackage[
    backend=biber,
    style=alphabetic,
    sorting=nyt
]{biblatex}

\addbibresource{references.bib}

% Theorem environments
\newtheorem{theorem}{Theorem}[section]
\newtheorem{lemma}[theorem]{Lemma}
\newtheorem{proposition}[theorem]{Proposition}
\newtheorem{corollary}[theorem]{Corollary}
\newtheorem{remark}[theorem]{Remark}

\theoremstyle{definition}
\newtheorem{definition}[theorem]{Definition}

\title{Computational Optimal Transport\\
Theoretical Notes and Proofs}
\date{\today}
\author{}

\begin{document}

\maketitle

\section{Discrete Optimal Transport}

We follow the framework of
\cite{peyre2019computational}
and \cite{santambrogio2015optimal}.

\subsection{Kantorovich Problem}

Let $n,m\geq 1$ and let$a\in\mathbb{R}_+^n$ and $b\in\mathbb{R}_+^m$ satisfy
\[\sum_{i=1}^{n}a_i=\sum_{j=1}^{m}b_j=M>0.\]

\begin{definition}[Transport Polytope]
The transport polytope associated with $(a,b)$ is 
\[ U(a,b) = \left\{P\in\mathbb{R}_+^{n\times m} : P\mathbf{1}_m = a, P^\top \mathbf{1}_n = b\right\}. \]
\end{definition}

Given a finite cost matrix $C\in\mathbb{R}^{n\times m}$, the discrete Kantorovich problem is
\begin{equation}
\label{eq:kantorovich}
\min_{P\in U(a,b)} \langle P, C\rangle = \sum_{i=1}^{n}\sum_{j=1}^{m}P_{ij}C_{ij}.
\end{equation}

\begin{lemma}
    \label{lem:transport-polytype}
    The transport polytope $U(a,b)$ is a non-empty, compact, convex polytope.
\end{lemma}

\begin{proof}
    We establish the assertions separately.

    \emph{Nonemptiness.}
    Define 
    \[ P^{(0)} = \frac{ab^\top}{M}. \]
    Since $a,b\geq 0$, we have $P^{(0)}\geq 0$. Moreover, we have
    \[ P^{(0)}\mathbf{1}_m = \frac{a(b^\top\mathbf{1}_m)}{M} = a, \]
    and
    \[ (P^{(0)})^\top\mathbf{1}_n = \frac{b(a^\top\mathbf{1}_n)}{M} = b. \]
    Hence, $P^{(0)}\in U(a,b)$, and $U(a,b)$ is non-empty.

    \emph{Convexity.}
    Let $P_1,P_2\in U(a,b)$ and $\lambda\in[0,1]$. Then, we set 
    \[ Q = \lambda P_1 + (1-\lambda)P_2. \] 
    Then $Q\geq 0$ and
    \[ Q\mathbf{1}_m = \lambda P_1\mathbf{1}_m + (1-\lambda)P_2\mathbf{1}_m = \lambda a + (1-\lambda)a = a, \]
    and similarly,
    \[ Q^\top\mathbf{1}_n = b. \]
    Hence, $Q\in U(a,b)$, and $U(a,b)$ is convex.

    \emph{Compactness.}
    For any $P\in U(a,b)$, nonegativity and the marginal constraints imply that
    \[ 0\leq P_{ij}\leq \min(a_i,b_j) \quad \forall i,j. \]
    Hence, $U(a,b)$ is bounded. 
    Since $U(a,b)$ is defined by a finite number of linear equalities and inequalities, 
    it is closed. Therefore, $U(a,b)$ is compact.

    Finally, since $U(a,b)$ is a non-empty, compact, convex set defined by linear inequalities, it is a convex polytope.
\end{proof}

\subsection{Extreme Points and Sparse Optimal Plans}

\begin{theorem}[Existence of Sparse Optimal Plans]
    \label{thm:sparse-optimal-plan}
    Let $a\in\mathbb{R}_+^n$ and $b\in\mathbb{R}_+^m$ have the same positive total mass,
    and let $C\in\mathbb{R}^{n\times m}$ be a cost matrix. 

    Then the Kantorovich problem \eqref{eq:kantorovich} admits an optimal solution 
    $P^\star\in U(a,b)$ which is an extreme point of the transport polytope $U(a,b)$.
    and satisfies
    \[\left|\operatorname{supp}(P^*)\right|\leq n+m-1.\]
\end{theorem}

\begin{proof}
    We first establish the rank of the linear constraints matrix.

    Let $p=\operatorname{vec}(P)\in\mathbb{R}^{nm}$ be the row-wise vectorization 
    of $P$. The marginal constraints can be expressed as a linear system 
    $Ap = \begin{pmatrix} a \\ b \end{pmatrix}$, where $A\in\mathbb{R}^{(n+m)\times nm}$.
    
    The column of $A$ associated with the variable $P_{ij}$ is 
    \[ A_{ij} = \begin{pmatrix} e_i^{(n)} \\ e_j^{(m)} \end{pmatrix}.\]

    We claim that the rank of $A$ is $n+m-1$.

    Indeed, let $(\lambda, \mu)\in\mathbb{R^n}\times\mathbb{R^m}$ satisfy
    \[(\lambda^\top, \mu^\top)A = 0.\]
    By the definition of $A$, this implies that $\lambda_i + \mu_j = 0$ for all $i,j$.
    Consequently, there exists $t\in\mathbb{R}$ such that $\lambda_i = t\mathbf{1}_n$
    and $\mu_j = -t\mathbf{1_m}$. Thus the left null space of $A$ is one-dimensional, 
    and the rank of $A$ is $n+m-1$.

    Now let $P\in U(a,b)$ be an extreme point and set 
    \[S=\operatorname{supp}(P)=\{(i,j):P_{ij}>0\}.\]

    Suppose for contradiction that $|S|\geq n+m$. 
    Since the columns of $A$ indexed by $S$ outnumber the rank of $A$, 
    they are linearly dependent.

    Hence there exists a non-zero $H\in\mathbb{R}^{n\times m}$ 
    such that \[\operatorname{supp}(H)\subseteq S\] and 
    \[H\mathbf{1}_m=0,\qquad H^\top\mathbf{1}_n=0.\]

    Since $P_{ij}>0$ for all $(i,j)\in S$, there exists $\epsilon>0$ such that
    \[ P\pm\epsilon H\geq 0. \] Moreover, 
    \[ (P\pm\epsilon H)\mathbf{1}_m = a, \qquad (P\pm\epsilon H)^\top\mathbf{1}_n = b.\]
    Therefore both $P+\epsilon H$ and $P-\epsilon H$ belong to $U(a,b)$.

    Since $H\neq 0$, we have $P+\epsilon H\neq P-\epsilon H$.
    But $P = \frac{1}{2}(P+\epsilon H) + \frac{1}{2}(P-\epsilon H)$, 
    which contradicts the assumption that $P$ is an extreme point.
    Therefore every extreme point $P\in U(a,b)$ satisfies 
    $|\operatorname{supp}(P)|\leq n+m-1$.

    Finally, by lemma \ref{lem:transport-polytype}, $U(a,b)$ is a non-empty, compact, convex polytope.
    A linear functional attains its minimum over a non-empty, compact, convex polytope at an extreme point.
    Thus there exists an optimal extreme point $P^\star\in U(a,b)$, and 
    \[|\operatorname{supp}(P^\star)|\leq n+m-1.\]
\end{proof}

\subsection{The Birkhoff--von Neumann Theorem}

We now consider the special case of the transport polytope where $n=m$ and $a=b=\mathbf{1}_n$. 
In this case, the transport polytope is known as the Birkhoff polytope.
\[\mathcal{B}_n = \left\{Q\in\mathbb{R}_+^{n\times n}:Q\mathbf{1}_n=\mathbf{1}_n,\;
Q^\top\mathbf{1}_n=\mathbf{1}_n\right\}.\]

The matrices in $\mathcal{B}_n$ are known as doubly stochastic matrices. By the change of variables $Q = nP$,
we have $U(a,b)=\frac{1}{n}\mathcal{B}_n$.

\begin{theorem}[Birkhoff--von Neumann]
\label{thm:birkhoff}
The extreme points of $\mathcal{B}_n$ are precisely the permutation matrices.
Consequently,
\[
\mathcal{B}_n
=\operatorname{conv}
\left\{\Pi_\pi:\pi\in S_n\right\},
\]
where $\Pi_\pi$ denotes the permutation matrix associated with $\pi\in S_n$.
\end{theorem}

\begin{proof}
We prove the characterization of extreme points in both directions.

\medskip
\noindent
\textit{Step 1: Every extreme point is a permutation matrix.}

Let $Q$ be an extreme point of $\mathcal{B}_n$.

Consider its bipartite support graph
$G(Q)=(X\sqcup Y,E)$, where
\[ X=\{x_1,\ldots,x_n\},\qquad Y=\{y_1,\ldots,y_n\},\]
and
\[(x_i,y_j)\in E\quad\Longleftrightarrow\quad Q_{ij}>0.\]

We first show that $G(Q)$ contains no cycles. Suppose, for contradiction, that
$G(Q)$ contains a cycle. Since $G(Q)$ is bipartite, every cycle has even length.

Assign alternating values $+1$ and $-1$ to the entries corresponding to the
edges of this cycle, and zero elsewhere. This defines a nonzero matrix
$H\in\mathbb{R}^{n\times n}$ satisfying
\[
H\mathbf{1}_n=0,\qquad H^\top\mathbf{1}_n=0.
\]

All entries of $Q$ on the cycleare strictly positive. Therefore, for sufficiently small
$\delta>0$, we have
\[Q+\delta H\geq 0,\qquad Q-\delta H\geq 0.\]

Moreover,
\[(Q\pm\delta H)\mathbf{1}_n=\mathbf{1}_n,\]
and
\[(Q\pm\delta H)^\top\mathbf{1}_n=\mathbf{1}_n.\]
Hence both $Q+\delta H$ and $Q-\delta H$ belong to $\mathcal{B}_n$.
Since $H\neq0$, these matrices are distinct, while
\[Q=\frac12(Q+\delta H)+\frac12(Q-\delta H).\]
This contradicts the extremality of $Q$. Thus $G(Q)$ is a forest.
Every connected component of this forest containing an edge has a leaf.

Suppose first that such a leafis a left vertex $x_i$, whose unique neighbor is $y_j$.
Then \[Q_{ik}=0\qquad\text{for all }k\neq j.\]
Since the $i$-th row sums to one,\[Q_{ij}=1.\]
As the $j$-th column also sums to oneand all entries are nonnegative,
\[Q_{kj}=0\qquad\text{for all }k\neq i.\]
Thus the edge $(x_i,y_j)$ forms an isolated matching pair.
If the leaf is a right vertex,the same conclusion followsby exchanging rows and columns.

Remove the corresponding row and column from $Q$. The remaining $(n-1)\times(n-1)$
matrix is again doubly stochastic, and its support graph is a forest.
Repeating this argument inductively shows that $Q$ has exactly one entry equal to one in each row
and column, with all other entries zero. Therefore $Q$ is a permutation matrix.

\medskip
\noindent
\textit{Step 2: Every permutation matrixis an extreme point.}

Let $\Pi$ be a permutation matrix. Then there exists $\pi\in S_n$ such that
\[\Pi_{ij}=
\begin{cases}
1, & j=\pi(i),\\
0, & j\neq\pi(i).
\end{cases}
\]

Suppose\[\Pi=\lambda Q+(1-\lambda)R,\] where $Q,R\in\mathcal{B}_n$ and $\lambda\in(0,1)$.
For every $j\neq\pi(i)$, we have $\Pi_{ij}=0$, hence
\[0=\lambda Q_{ij}+(1-\lambda)R_{ij}.\]
Since $Q_{ij},R_{ij}\geq0$and both coefficients are strictlypositive, it follows that
\[Q_{ij}=R_{ij}=0\qquad(j\neq\pi(i)).\]
Since each row of $Q$ and $R$ sums to one, we obtain
\[Q_{i,\pi(i)}=R_{i,\pi(i)}=1.\]
Consequently,
\[Q=R=\Pi.\]
Therefore $\Pi$ is an extreme pointof $\mathcal{B}_n$.

We conclude that
\[\operatorname{Ext}(\mathcal{B}_n)=\{\Pi_\pi:\pi\in S_n\}.\]

Finally, every nonempty compactpolytope is the convex hull ofits extreme points. Hence
\[\mathcal{B}_n=\operatorname{conv}\{\Pi_\pi:\pi\in S_n\}.\]
\end{proof}

\begin{corollary}[Permutation-type optimal plan]
\label{cor:permutation-optimal}
Suppose $n=m$ and\[a=b=\frac1n\mathbf{1}_n.\]
Then the Kantorovich problem admitsan optimal transport plan of the form
\[P^*=\frac1n\Pi_{\pi^*},\]where $\Pi_{\pi^*}$ is a permutation matrix.
\end{corollary}

\begin{proof}
Since the Kantorovich objective \[P\longmapsto\langle C,P\rangle\]
is linear and $U(a,b)$ is a nonempty compact polytope,
an optimal solution exists at an extreme point.

By Theorem~\ref{thm:birkhoff}, every extreme point of
\[U(a,b)=\frac1n\mathcal{B}_n\] is of the form
\[P=\frac1n\Pi_\pi \] for some $\pi\in S_n$.

Therefore, there exists an optimal permutation $\pi^*$ satisfying
\[\pi^*\in\operatorname*{arg\,min}_{\pi\in S_n}\sum_{i=1}^n C_{i,\pi(i)},\]
and
\[P^*=\frac1n\Pi_{\pi^*}.\]
\end{proof}

\section{Kantorovich Duality}
\subsection{Derivation of the Dual Problem}

Recall the discrete Kantorovich problem
\begin{equation}
\label{eq:primal-ot}
\begin{aligned}
\min_{P\in\mathbb{R}_+^{n\times m}}
\quad & \langle C,P\rangle\\
\text{subject to}\quad
& P\mathbf{1}_m=a,\\
& P^\top\mathbf{1}_n=b,
\end{aligned}
\end{equation}
where $a\in\mathbb{R}_+^n$ and $b\in\mathbb{R}_+^m$ have equal positive
total mass and $C\in\mathbb{R}^{n\times m}$ is a finite cost matrix.

Introduce Lagrange multipliers$f\in\mathbb{R}^n$ and $g\in\mathbb{R}^m$ for the two sets of marginal constraints.
The Lagrangian is
\[
\begin{aligned}
L(P,f,g)={}&\langle C,P\rangle+\langle f,a-P\mathbf{1}_m\rangle +\langle g,b-P^\top\mathbf{1}_n\rangle.
\end{aligned}
\]
Rearranging the terms gives
\[
L(P,f,g)=\langle a,f\rangle+\langle b,g\rangle +\sum_{i=1}^n\sum_{j=1}^m (C_{ij}-f_i-g_j)P_{ij}.
\]

The dual function is defined by \[q(f,g)=\inf_{P\geq0}L(P,f,g).\]
Therefore, the Kantorovich dual problem is
\begin{equation}
\label{eq:dual-ot}
\begin{aligned}
\max_{f\in\mathbb{R}^n,\,
      g\in\mathbb{R}^m}
\quad&
\langle a,f\rangle+\langle b,g\rangle\\
\text{subject to}\quad&
f_i+g_j\leq C_{ij}
\quad\text{for all }i,j.
\end{aligned}
\end{equation}

\subsection{Weak Duality}

\begin{proposition}[Weak duality]
\label{prop:weak-duality}
For every primal feasible
$P\in U(a,b)$ and every dual feasible
pair $(f,g)$,
\[
\langle a,f\rangle+\langle b,g\rangle
\leq \langle C,P\rangle.
\]
\end{proposition}

\begin{proof}
Since $(f,g)$ is dual feasible,
\[
f_i+g_j\leq C_{ij}
\qquad\text{for all }i,j.
\]
Multiplying by $P_{ij}\geq0$ and summing over all indices yields
\[
\sum_{i,j}(f_i+g_j)P_{ij}\leq\sum_{i,j}C_{ij}P_{ij}.
\]
Using the marginal constraints,
\[
\begin{aligned}
\sum_{i,j}(f_i+g_j)P_{ij}
&=\sum_i f_i\sum_jP_{ij}+\sum_jg_j\sum_iP_{ij}\\
&=\sum_i a_if_i+\sum_jb_jg_j.
\end{aligned}
\]
Consequently,
\[
\langle a,f\rangle+\langle b,g\rangle\leq\langle C,P\rangle.
\]
\end{proof}

\subsection{Strong Duality}

\begin{theorem}[Strong Kantorovich duality]
\label{thm:strong-duality}
Let $a\in\mathbb{R}_+^n$ and$b\in\mathbb{R}_+^m$ have the same
positive total mass, and let$C\in\mathbb{R}^{n\times m}$.
Then the primal and dual problems\eqref{eq:primal-ot} and \eqref{eq:dual-ot} both attain
their optimal values, and
\[
\min_{P\in U(a,b)}\langle C,P\rangle=
\max_{\substack{f\in\mathbb{R}^n,\,
g\in\mathbb{R}^m\\
f_i+g_j\leq C_{ij}}}
\left(
\langle a,f\rangle+\langle b,g\rangle
\right).
\]
\end{theorem}

\begin{proof}
By the compactness and nonemptinessof $U(a,b)$ established previously,
the primal problem admits an optimal solution $P^*$.

Let $p^*=\operatorname{vec}(P^*)$,$c=\operatorname{vec}(C)$, and write
the marginal constraints as
\[
Ap=d,\qquad d=
\begin{pmatrix}
a\\b
\end{pmatrix}.
\]
Define the index sets
\[ I=\{k:p_k^*>0\}, \qquad Z=\{k:p_k^*=0\}. \]
Consider the cone of feasible directions
\[
\mathcal{D}=\left\{h\in\mathbb{R}^{nm}:Ah=0,\; h_k\geq0\ \text{for }k\in Z\right\}.
\]
For any $h\in\mathcal{D}$,there exists a sufficiently small
$t>0$ such that
\[
p^*+th\geq0,
\qquad
A(p^*+th)=d.
\]

Since $p^*$ is optimal,
\[
c^\top(p^*+th)\geq c^\top p^*.
\]
Hence
\[
c^\top h\geq0 \qquad\text{for all }h\in\mathcal{D}.
\]
By the dual-cone form of Farkas' lemma, there exist
$y\in\mathbb{R}^{n+m}$ and $s\in\mathbb{R}_+^{nm}$ such that
\[c=A^\top y+s,\] with
\[s_k=0\qquad\text{for every }k\in I.\]

Write
\[
y=
\begin{pmatrix}
f^*\\g^*
\end{pmatrix}.
\]
By the structure of the marginalconstraint matrix $A$,
\[
(A^\top y)_{ij}=f_i^*+g_j^*.
\]
It follows that
\[
C_{ij}=f_i^*+g_j^*+s_{ij},\qquad s_{ij}\geq0.
\]
Therefore
\[
f_i^*+g_j^*\leq C_{ij}\qquad\text{for all }i,j,
\]
so $(f^*,g^*)$ is dual feasible.

Moreover, the definition of $I$ and the condition on $s$ imply
\[
s^\top p^*=0.
\]
Thus
\[
\begin{aligned}
c^\top p^*
&=(A^\top y+s)^\top p^*\\
&=y^\top Ap^*+s^\top p^*\\
&=y^\top d\\
&=\langle a,f^*\rangle
+\langle b,g^*\rangle.
\end{aligned}
\]
\end{proof}

\subsection{Complementary Slackness}

\begin{proposition}[Complementary slackness]
\label{prop:complementary-slackness}
Let $P^*\in U(a,b)$ be primal optimal, and let $(f^*,g^*)$ be dual optimal.
Then\[P_{ij}^*(C_{ij}-f_i^*-g_j^*)=0\qquad\text{for all }i,j.\]
Conversely, any primal feasible$P$ and dual feasible $(f,g)$satisfying these identitiesare both optimal.
\end{proposition}

\begin{proof}
By strong duality,
\[
\langle C,P^*\rangle=\langle a,f^*\rangle+\langle b,g^*\rangle.
\]
Using the marginal constraints,we obtain
\[
\sum_{i,j}P_{ij}^*(C_{ij}-f_i^*-g_j^*)=0.\]
Each term in this sum is nonnegative,because $P^*$ is primal feasible and $(f^*,g^*)$ is dual feasible.

Therefore every term must vanish:\[P_{ij}^*(C_{ij}-f_i^*-g_j^*)=0.\]
Conversely, suppose $P$ and $(f,g)$are feasible and satisfy these complementary slackness identities.
Then
\[
\begin{aligned}
\langle C,P\rangle
&=\sum_{i,j}(f_i+g_j)P_{ij}\\
&=\langle a,f\rangle+\langle b,g\rangle.
\end{aligned}
\]
By weak duality, equality ofthe primal and dual objective valuesimplies that both solutions are optimal.
\end{proof}


\end{document}
